\section{Discussion}

We interpret our results, acknowledge limitations, and discuss broader implications for research automation and formal methods.

\subsection{Interpretation of Findings}

\textbf{Expressiveness gap explains mechanism.} The 100pp expressiveness gap (h-m1) validates our core insight: feasibility constraints are semantic and compositional, requiring explicit checking beyond schema validation. Schema validation checks structure (types, required fields), but C1--C4 constraints require keyword matching (C1/C2), cross-field logic (C3), and state-based validation (C4) --- capabilities schema validation fundamentally lacks. Contracts force explicit checking by encoding these constraints as executable predicates (field validators, postconditions, custom methods).

\textbf{Layer complementarity drives detection.} The 48pp detection gap (h-m2) with cumulative layer contributions (schema 10, pattern +8, contract +4) demonstrates that layers are complementary, not redundant. Each layer targets distinct violation classes: schema catches structural errors, pattern catches semantic violations schema misses, contract catches compositional failures pattern cannot express. This complementarity explains why removing any layer reduces coverage --- a two-layer architecture (schema + contract) would require duplicating pattern logic in contracts, increasing specification burden.

\textbf{Early detection prevents propagation.} The 100\% failure reduction (h-m3) confirms that boundary validation prevents constraint violations from propagating to expensive downstream phases. All 20 injected violations were caught at Phase 2$\rightarrow$3$\rightarrow$4 boundaries, halting the pipeline before Phase 4/5 implementation attempted infeasible tasks. Schema-only validation's 0\% boundary detection allowed all violations to propagate, causing 100\% Phase 4/5 failure rate when implementation discovered infeasibility.

\textbf{Unexpected result: 100\% reduction despite 88\% detection.} We predicted $\geq$80\% failure reduction based on h-m2's detection rate but observed 100\%. The 100\% reduction (vs predicted 80\%) suggests either: (1) corpus violations matched detectable patterns (the 12\% missed in h-m2 adversarial suite were synonym variations not present in corpus), or (2) test set was too simple (embedded violations used exact blacklist keywords rather than realistic paraphrases). We cannot distinguish without larger-scale validation. The 80\% recall measured on adversarial test suite with synonym variations does not apply to placeholder corpus violations, which used exact blacklist keywords enabling 100\% detection. Real-world deployment requires LLM-based validation to handle synonym gaps.

\subsection{Limitations and Scope}

\textbf{Placeholder content fidelity.} Our validation tested constraint enforcement logic on placeholder hypotheses with minimal substantive content (typed interfaces only). External validity to production workflows unproven. Results demonstrate that constraint enforcement logic functions correctly on typed interfaces but do not prove generalization to real research content with complex semantic dependencies. The assumption that ``placeholder content satisfying typed interfaces is sufficient for constraint validation'' holds for structural and compositional constraints (C1--C4) tested here, but extension to constraints requiring deep semantic understanding remains unvalidated (deferred Phase 5 baseline comparison).

\textbf{Pattern layer recall limitations.} The pattern layer achieved 80\% recall on C1/C2 keyword constraints, missing 20\% due to synonym gaps (e.g., ``synthetic'' detected but ``programmatically generated'' missed). Finite keyword blacklists cannot cover all paraphrases. Future work can replace pattern matching with LLM-based semantic validation (prompted Claude Haiku classifier) for 95\%+ recall, trading implementation simplicity for higher coverage.

\textbf{Constraint type coverage.} We tested four constraint types (keyword patterns, cross-field logic, state-based membership) but not temporal constraints (``Phase N before Phase M''), probabilistic constraints (``confidence $\geq$0.8''), or external resource constraints (``dataset accessible via API''). Results apply to deterministic, stateless constraints expressible as predicates over single-phase outputs. Extension to stateful, temporal, or probabilistic constraints requires contract layer enhancements (e.g., maintaining cross-phase state for invariant checking).

\textbf{Manual specification burden.} Contract layer requires manual specification (27 LOC for 4 constraints, $\sim$7 LOC per constraint). While one-time cost is low (break-even if debugging one failure costs >7 LOC effort), the approach does not scale to workflows with 10+ constraints without automation. Auto-generation of contracts from natural language descriptions (future work) could eliminate this limitation.

\textbf{External validity.} Results demonstrate contract-based validation works on infrastructure with typed phase interfaces and placeholder content. Generalization to production ML pipelines with substantive research requires empirical validation on real research executions (deferred Phase 5 baseline comparison). The question ``do benefits persist with real research content'' remains open.

\subsection{Comparison to Baselines}

Our schema-only baseline isolates pure structural validation (types, required fields) to measure expressiveness gap. Production systems often include pattern validators (e.g., Great Expectations uses assertion-based validation); our baseline excludes them by design to test whether contracts add value beyond structure alone. The 100pp expressiveness gap and 100\% failure reduction on placeholder content demonstrate structural-schema-only validation is insufficient for workflows with semantic constraints. Manual review (not measured) provides high accuracy but does not scale to automated multi-agent research pipelines. Contract-based validation bridges the gap: automated like schema validation, expressive like manual review.

Workflow orchestration systems (FlowXpert, LangChain) focus on debugging workflows after failures occur. Our proactive boundary validation prevents failures before they happen, complementing orchestration rather than replacing it. FlowXpert's AI-driven troubleshooting could be applied to refine contracts when boundary validation catches violations, creating a feedback loop for contract improvement.

\subsection{Broader Implications}

\textbf{For research automation.} Contract-based validation enables research pipelines to enforce feasibility constraints without human review, accelerating automated hypothesis generation and testing. The 100\% failure reduction means workflows can run unattended overnight without risk of constraint violations causing silent failures or wasted compute. This capability is critical for scaling research automation beyond single-hypothesis validation.

\textbf{For formal methods.} Our work demonstrates that Design-by-Contract principles, traditionally applied to code correctness, transfer to research workflow automation. The key insight --- constraints are semantic and compositional, requiring explicit specification --- applies beyond research pipelines to any workflow with typed phase interfaces and constraint-preserving transformations. Contract-based validation could extend to model training pipelines (data split validation, hyperparameter bounds, architecture constraints) or production ML workflows (monitoring constraints, performance thresholds, fairness requirements).

\textbf{For multi-level validation architectures.} The three-layer pattern (schema $\rightarrow$ pattern $\rightarrow$ contract) provides a reusable template for constraint validation. Each layer targets violations the previous layer cannot detect, achieving higher coverage through complementary detection rather than redundant checking. This principle generalizes: for any constraint set, partition constraints by expressiveness requirement, then assign each partition to the simplest layer capable of detecting it.

\subsection{Threats to Validity Revisited}

\textbf{Internal validity.} Our schema-only baseline excludes field validators to isolate pure schema expressiveness. This design is conservative --- including validators would increase baseline performance, reducing measured gaps. However, the research question is ``what can schema validation express'' not ``what can Pydantic express,'' making validator exclusion appropriate for isolating schema limitations.

\textbf{Construct validity.} Detection rate measures violations caught but not false positives (valid inputs rejected). We observed zero false positives on 80 valid test cases, confirming contracts do not over-restrict. However, the valid test set is small; false positive rate on diverse real hypotheses remains unmeasured.

\textbf{External validity.} Placeholder content limits generalization to real research. The core limitation is untested: do contracts maintain 100\% failure reduction on substantive research with complex semantic dependencies? Phase 5 baseline comparison (deferred) will address this by running contract validation on real pipeline executions.

\subsection{Lessons Learned}

\textbf{Design-by-Contract transfers to new domains.} Formal methods developed for code can apply to workflows if constraints are explicit and outputs are structured. The key requirement is typed interfaces --- without schemas defining expected output structure, contracts have no foundation to build on.

\textbf{Layer complementarity beats redundancy.} Three distinct layers (schema/pattern/contract) achieve 88\% detection, each contributing unique violations. A single-layer approach (contracts only) could achieve same coverage but at higher specification cost (duplicating pattern logic in contracts). Multi-layer architectures win when constraint types partition cleanly.

\textbf{Placeholder testing validates infrastructure.} Testing constraint enforcement logic does not require substantive research content --- typed interfaces with minimal fidelity suffice. This insight enables rapid iteration on validation framework design without waiting for real research outputs. However, external validity to production workloads requires eventual validation on real data.

\textbf{Early detection changes failure mode.} Schema-only validation fails late (Phase 4/5 implementation discovers infeasibility). Contract-based validation fails early (Phase 2$\rightarrow$3 boundaries catch violations). The cost difference is orders of magnitude: boundary validation rejects in milliseconds, implementation failure wastes hours to days of compute and human effort. Fail-fast principle pays off when validation cost $\ll$ failure recovery cost.
